\HMTNarrativeHeading{La profondeur d’un nombre}

La génération coinductive organise la détermination d’un nombre au moyen de
prolongements compatibles. Une étape produit un préfixe et conserve les
conditions pour le continuer. La profondeur peut augmenter tout en maintenant
la concordance de toutes les lectures précédentes. La valeur apparaît lors de la
réalisation de cette compatibilité ; l’histoire conserve, en outre, la manière dont
la précision a été obtenue.

Les régions des coordonnées circulaire, d’or et exponentielle proviennent du
même développement APP--TRIT--TPK. Leurs lecteurs remplissent des fonctions
différenciées et partagent la structure qui permet de les relier.
L’extension triadique d’Euler explicite les réalisations des régimes ;
le calendrier aréal--volumétrique détermine l’auto-échelle ; les émissions
régionales et leurs frontières conservent les conditions d’évaluation.
La succession de ces étapes explicite l’ordre de construction des
coordonnées qui interviennent ensuite dans alpha.

Le calendrier des vacances permet de comprendre une particularité du
raffinement. L’histoire avance à chaque transition, tandis que la capacité
d’une lecture peut rester fixe. Le mot qui enregistre ces
transitions distingue la capacité acquise de la vacance conservée.
Son bilan compte les deux contributions. La régularité du calendrier et
l’apériodicité de son itinéraire décrivent des aspects liés d’une
même évolution.

La monodromie prolonge cette distinction. Revenir à une phase signifie reconnaître une position du cycle ; récupérer une histoire signifie restituer également ses tours, son orientation et ses registres. Le miroir régional conserve l'axe et l'agrégat des deux coordonnées qu'il échange, tandis qu'il inverse leur différence. Cette donnée permet de reconstruire la répartition. La figure nonadique situe les phases et leurs opérations ; ses marques sont des positions du calendrier, avec une sémantique différente de la valeur réelle d'une constante.

La cinquième frontière, les régions ultérieures et le retour s’expliquent
donc au moyen de leurs actions sur l’état. Les figures régionales
montrent simultanément la coïncidence d’une lecture et la différenciation de
ses antécédents. Cette différenciation transmet à la réalisation
exceptionnelle l’information incidentielle qu’elle utilise. La génération
numérique poursuit ainsi la double projection posée au début : chaque
lecture obtient sa valeur d’une structure dont l’organisation demeure
disponible pour les autres réalisations.

